\documentclass[12pt]{article}

\input{../../../_assets/preambulo}

\makeatletter
\let\ejercicio\undefined
\let\endejercicio\undefined
\let\c@ejercicio\undefined % Anula el contador interno

\expandafter\let\csname ejercicio*\endcsname\undefined
\expandafter\let\csname endejercicio*\endcsname\undefined
\makeatother

% 2. Creamos la macro del nombre directamente con el idioma deseado
\newcommand{\nameEjercicio}{Exercise}

% 3. Volvemos a declarar los entornos ahora que la vía está libre
\newtheorem*{ejercicio*}{\nameEjercicio}             % No numerado
\newtheorem{ejercicio}{\nameEjercicio} [section]     % Se resetea en cada section

\renewcommand{\theejercicio}{%
  \ifnum\value{section}=0 
    \arabic{ejercicio}% Solo mostrar el número de ejercicio
  \else
    \thesection.\arabic{ejercicio}% Mostrar número de sección y número de ejercicio
  \fi
}


\begin{document}

    % 1. Foto de fondo
    % 2. Título
    % 3. Encabezado Izquierdo
    % 4. Color de fondo
    % 5. Coord x del titulo
    % 6. Coord y del titulo
    % 7. Fecha

    
	\input{../../../_assets/portada}
	\portadaExamenFotoDif{../../sapienza.jpg}{Statistica \\Matematica\\Exam I}{Statistica Matematica. Exam I}{MidnightBlue}{-8}{28}{2025-2026}{Joaquín Avilés de la Fuente}

    \begin{description}
        \item[Asignatura] Statistica Matematica.
        \item[Curso Académico] 2025-2026.
        \item[Grado] Grado en Matemáticas.
        \item[Grupo] Grupo Único.
        \item[Profesor] Antonio Agresti.
        \item[Descripción] Examen final escrito de Statistica Matematica.
        \item[Fecha] 21 de enero de 2026.
        \item[Duración] 120 min.
    
    \end{description}
    \newpage

    % ------------------------------------
    
    \begin{ejercicio}[9 points]
        Let $X_1, \dots, X_n$ be a sample of independent and identically distributed random variables according to a $\text{Poisson}(\lambda)$ distribution, with $\lambda$ being an unknown parameter. Assume an exponential distribution with mean $1/2$ as the prior distribution for $\lambda$.
        \begin{itemize}
            \item (4 points) Determine the posterior distribution for the parameter $\lambda$ in terms of $n$ and the empirical mean of the sample.
            \item (4 points) Determine the Bayesian estimator for $\lambda$. Show that it is asymptotically unbiased and consistent.
            \item (1 point) Evaluate the bias. Does there exist a value of $\lambda$ for which the estimator is exactly unbiased?
        \end{itemize}
    Recall that 
    $$\Gamma(z) = \int_0^\infty dt \, t^{z-1} e^{-t}$$ 
    and $\Gamma(n+1) = n!$.

    \end{ejercicio}

    \begin{ejercicio}[6 points]
        Let $X_1, \dots, X_n$ be a sample of independent and identically distributed random variables with density:
        \[
        f_\theta(x) = \theta \, x^{\theta-1}, \qquad 0 < x < 1, \; \theta > 0.
        \]
        \begin{itemize}
            \item (3 points) Determine the maximum likelihood estimator $\hat{\theta}_{MLE}$.
            \item (3 points) Calculate the Fisher information for the sample.
        \end{itemize}
    \end{ejercicio}

    \begin{ejercicio}[8 points]
        Let $(X_n)_{n \ge 1}$ be a sequence of i.i.d. random variables. Consider: $Y_n = \left(\prod\limits_{i=1}^n X_i \right)^{1/n}$
        \begin{itemize}
            \item Determine $\lim\limits_{n \to \infty} Y_n$ almost surely in the following cases:
            \begin{itemize}
                \item (3 points) $X_n$ is uniformly distributed on $(0,1)$.
                \item (2 points) $X_n$ is distributed as $e^Z$ where $P(Z = k) = \dfrac{1}{2^k}$ for $k \in \mathbb{N}$ (Geometric).
            \end{itemize}
            \item (3 points) In the case where $X_n$ is uniformly distributed on $(0,1)$, determine the limit in distribution of:
            \[
            \sqrt{n}(Y_n - e^{-1}).
            \]
        \end{itemize}
        \textit{\small Hint for the last point: Use the fact that if a sequence of random variables $(W_n)_n$ satisfies $$\sqrt{n}(W_n - \theta) \xrightarrow{d} N(0, \sigma^2)$$ in distribution, then $$\sqrt{n}(f(W_n) - f(\theta)) \xrightarrow{d} N(0, \sigma^2 [f'(\theta)]^2)$$ in distribution, for any continuously differentiable function $f$.}
    \end{ejercicio}

    \begin{ejercicio}[7 points]
        Let $(\mathbb{R}_+^2, \mathcal{B}(\mathbb{R}_+^2), \text{Unif}_{[0,a]} \otimes \text{Unif}_{[0,b]}, (a,b) \in \mathbb{R}_+^2)$ be a statistical model.
        \begin{itemize}
            \item Identify a Pivot distribution for the parameter $\theta = (a,b)$ and determine a (non-trivial) confidence region of level $\gamma \in (0,1)$.
            \item Let $\tau(a,b) = (a+b, a-b)$ be a characteristic of the model. Determine a Pivot distribution for $\tau$ and determine a (non-trivial) confidence region of level $\gamma \in (0,1)$.
        \end{itemize}
    \end{ejercicio}

\end{document}